\section{Introduction}

Coarse index theory relates the large-scale geometry of a space to the
K-theory of an algebra of controlled operators. Its assembly map sends
geometric K-homology classes to their analytic indices. The injectivity
question asks whether that index retains the information in a class after
one has passed to large scales. Roe developed this viewpoint and its
Novikov-type formulation \cite{Roe93}; Higson and Roe clarified
the role of bounded geometry and coarsened K-homology \cite{HR95}.
Here we construct a bounded-geometry space for which ordinary reduced
coarse assembly loses a class of infinite order.

A metric space $X$ is uniformly discrete if distinct points have distance
bounded below by a positive constant. It has bounded geometry if, for every
$R<\infty$, the cardinalities of its radius-$R$ balls have a common finite
bound. For such a space, let $P_r(X)$ be the Rips complex: its simplices are the
finite subsets of $X$ of diameter at most $r$. We use locally finite analytic
K-homology and its coarsening,
\[
 K_i^{\lf}(Y)=\KK_i(C_0(Y),\C),\qquad
 KX_i(X)=\varinjlim_{r\to\infty}K_i^{\lf}(P_r(X)).
\]
Here and throughout, $i$ is taken modulo two. The reduced, locally compact
Roe algebra $C^*(X)$ is the operator-norm closure of the locally compact
finite-propagation operators on $\ell^2(X)\otimes\HH_0$, where $\HH_0$ is an
infinite-dimensional separable Hilbert space. Section~\ref{sec:conventions}
specifies these conventions and the ordinary coarse assembly map
\[
 \mu_X:KX_i(X)\longrightarrow K_i(C^*(X)).
\]
The coarse Baum--Connes conjecture asks whether this map is an isomorphism
\cite[Conjecture~6.6]{HR95}. The rational coarse Novikov assertion is that
$\mu_X\otimes\id_{\Q}$ is injective. Passing through all Rips scales is
essential: a cycle at one scale can bound at a larger scale and therefore
represent zero in $KX_*$. An obstruction to injectivity must survive every
such coarsening while its analytic index vanishes.

Yu proved coarse assembly isomorphism for bounded-geometry spaces of finite
asymptotic dimension \cite[Theorem~7.1]{Yu98}, and subsequently for those admitting a
coarse embedding into Hilbert space \cite{Yu00}. Skandalis, Tu and Yu placed
the latter result in a groupoid framework \cite{STY02}. Further positive
injectivity results include coarse embeddings into uniformly convex Banach
spaces \cite[Theorem~1.1]{KY06} and fibred coarse embeddings into Hilbert space
\cite[Theorem~34]{FS14}. These theorems impose geometric hypotheses beyond bounded
geometry. Our result shows that bounded geometry alone does not suffice.

We prove the following integral statement.

\begin{theorem}\label{thm:main}
There is a uniformly discrete bounded-geometry metric space $X$, which is a
coarse disjoint union of finite connected graphs of uniformly bounded degree,
and an infinite-order element $\alpha\in KX_1(X)$ such that
\[
 \mu_X(\alpha)=0\quad\text{in }K_1(C^*(X)).
\]
Consequently, ordinary coarse assembly for $X$ fails to be injective both
integrally and after tensoring its domain and target with $\Q$.
\end{theorem}

\subsection*{Predecessors and the scope of the obstruction}
Yu already constructed a non-bounded-geometry example with a nontrivial
coarse K-homology class of zero index, using spheres of increasing dimension
and size \cite[Section~8, Proposition~8.1]{Yu98}. Dranishnikov, Ferry and
Weinberger constructed integral noninjectivity examples for a map from
locally finite K-homology to Roe-algebra K-theory, also allowing spaces
without bounded geometry \cite[Introduction, pp.~920--921]{DFW}.
Their Remark~6.2 distinguishes those examples from the Rips-coarsened
formulation: the relevant classes die under coarsening. Theorem~\ref{thm:main}
retains bounded geometry and detects its class in the actual direct limit.

Higson, Lafforgue and Skandalis constructed injectivity and surjectivity
obstructions for several groupoid assembly maps \cite{HLS}. Their Hausdorff
groupoid example in Section~5 combines a degree-one Bott class with a
Kazhdan projection to produce a kernel class. This is a close methodological
antecedent of the winding and constant-vector averaging used here.
Their coarse-space construction in Section~6 instead establishes
non-surjectivity for specific congruence expander unions
\cite[Section~6, Proposition~11]{HLS}. These conclusions must be separated:
a groupoid kernel example does not itself provide the bounded-geometry
space and ordinary Rips-domain kernel in Theorem~\ref{thm:main}.

Expander phenomena also coexist with injectivity. Willett and Yu proved
ordinary coarse injectivity for coarse disjoint unions of finite graphs
whose girth tends to infinity; with an additional weak-expansion hypothesis,
surjectivity fails \cite[Theorem~1.5]{WY12I}. The fibers in our construction
have short cycles, so persistent boundary winding does not imply that
large-girth hypothesis. More recent isomorphism results for certain relative
expanders likewise depend on specific embedding hypotheses for the kernel
and quotient families \cite[Theorem~1.1]{DWY23}. The organization of our
vertices into finite fibers does not assert such an embedding or a finite-group
extension. Throughout, our target is the reduced, locally compact Roe algebra;
maximal completions and coefficient-free group assembly are different
questions.
The independent companion article~\cite[Theorem~1.1]{OpenAITorsionFreeBC2026}
concerns group assembly, which for a torsion-free group maps the compactly supported K-homology
of its classifying space to the K-theory of its reduced group algebra. It
constructs a finitely generated torsion-free group for which this
coefficient-free map is not rationally injective.

\subsection*{How the construction works}
Begin with finite discs $D_j$ of radii $j$ in the triangular lattice, and let $C_j$ be their
boundary cycles. Above each vertex of $D_j$ place a finite congruence quotient
of $\SL_3(\Z)$. The quotient is trivial on $C_j$ and grows with depth inside
the disc. Cayley edges connect vertices within a fiber; reduction maps and
twisted matchings connect neighboring fibers. The resulting finite graph is
denoted by $X_j$. Thus $C_j$ is also a cycle in $X_j$.

Let $B$, $D$ and $X$ be the separated unions of the $C_j$, $D_j$ and $X_j$,
respectively. The family of oriented boundary classes gives an element of
$K_1^{\lf}(P_1(B))$, and its image defines $\alpha\in KX_1(X)$. Two features
of the fibers have different roles.

First, the twisting prevents short loops in $X_j$ from winding around the
center once $j$ is sufficiently large. At any prescribed Rips scale this
produces an integral cocycle whose restriction to $C_j$ has winding one.
The associated K-theory pairing detects every nonzero integer multiple of
$\alpha$ at that scale. Sections~\ref{sec:geometry} and
\ref{sec:persistence} prove that $\alpha$ has infinite order in the direct
limit.

Second, property~(T) uniformly approximates the projections onto constant
vectors in all the finite fibers. Averaging into those vectors defines an
isometry $V:\ell^2(D)\otimes\HH_0\to\ell^2(X)\otimes\HH_0$.
The right translations used for twisting fix the constant vectors, so they
do not interfere with this averaging.
Section~\ref{sec:transfer} proves that conjugation by $V$ maps $C^*(D)$ into
$C^*(X)$, despite the unbounded sizes of the fibers. On the boundary fibers
this isometry is the ordinary inclusion. Since the boundary classes bound
the filled discs, naturality of assembly then gives $\mu_X(\alpha)=0$ in
Section~\ref{sec:vanishing}.

The transfer argument isolates a useful principle: uniform finite-propagation
approximation of constant-fiber projections, together with controlled
operators between adjacent fibers, can produce a Roe-algebra homomorphism
without a corresponding coarse map. This principle accounts for the
difference between the persistent winding in $X$ and the vanishing of its
index through $D$.


Section~\ref{sec:conventions} fixes the analytic K-homology, Rips metrics,
and index naturality used by both branches. Section~\ref{sec:geometry}
constructs the common graph family and its short-loop estimate.
Section~\ref{sec:persistence} then proves persistence, independently of
averaging. Section~\ref{sec:transfer} proves the operator transfer, and
Section~\ref{sec:vanishing} combines its exact boundary restriction with
the filled-disc argument to annihilate the same class.
